\chapter{Sesión 4: Canales, suscripciones y publicaciones}
\label{ch:sesion4}

\section{Objetivo de la sesión}

Implementar los endpoints REST del dominio principal de NotiUAM:

\begin{itemize}[leftmargin=*]
  \item CRUD de canales temáticos (\texttt{/api/channels}).
  \item Suscripción y desuscripción idempotentes (\texttt{/api/channels/\{id\}/subscribe}).
  \item Publicación y listado paginado de avisos (\texttt{/api/channels/\{id\}/posts}).
  \item Control de acceso basado en el usuario autenticado (dueño del canal, autor del post, rol admin).
\end{itemize}

\section{Contexto}

En la Sesión~\ref{ch:sesion3} quedó listo el sistema de autenticación JWT. Sin embargo, el backend solo exponía \texttt{/api/auth/**}; no existía ningún endpoint de negocio. Tampoco había una forma limpia de obtener el usuario autenticado desde los servicios, ni excepciones específicas del dominio para mapear a códigos HTTP (404, 403).

\section{Desarrollo}

\subsection{Estructura creada}

\begin{lstlisting}[language=bash]
src/main/java/mx/uam/notiuam/
|- config/
|  `- CurrentUser.java                 (helper de usuario autenticado)
|- dto/channel/
|  |- ChannelCreateRequest.java
|  |- ChannelUpdateRequest.java
|  |- ChannelResponse.java
|  `- ChannelSummaryResponse.java
|- dto/post/
|  |- PostCreateRequest.java
|  `- PostResponse.java
|- service/
|  |- ChannelService.java
|  |- SubscriptionService.java
|  `- PostService.java
|- controller/
|  `- ChannelController.java
`- exception/
   |- NotFoundException.java
   `- ForbiddenException.java
\end{lstlisting}

\subsection{Decisiones técnicas}

\begin{itemize}[leftmargin=*]
  \item \textbf{Separación controller / service.} Los controllers no acceden directamente a los repositorios; toda la lógica vive en la capa de servicio.
  \item \textbf{UUID como identificador público.} A diferencia de un \texttt{Long autoincremental}, los UUID evitan que se pueda enumerar recursos (\texttt{/api/channels/1}, \texttt{/2}, ...).
  \item \textbf{Suscripción idempotente.} Llamar dos veces al endpoint de suscribirse no devuelve error; el segundo intento es un no-op. Esto simplifica el cliente móvil y evita problemas por reintentos de red.
  \item \textbf{Control de acceso en servicio, no en controller.} El \texttt{ChannelService} sabe quién es el dueño del canal y quién es admin; el controller solo delega.
  \item \textbf{Paginación en posts.} Los listados de publicaciones usan \texttt{Pageable} de Spring Data con un tope de 50 elementos por página para evitar respuestas gigantes.
\end{itemize}

\subsection{Componentes clave}

\subsubsection{\texttt{CurrentUser}}

Helper inyectable que lee el \texttt{SecurityContext} y devuelve el principal tipado. Antes, cada controller tendría que hacer un \texttt{cast} manual; ahora basta con inyectar y llamar a \texttt{currentUser.requireId()}.

\begin{lstlisting}[language=java]
public AuthenticatedUser require() {
    Authentication auth = SecurityContextHolder.getContext().getAuthentication();
    if (auth == null || !(auth.getPrincipal() instanceof AuthenticatedUser u)) {
        throw new AccessDeniedException("No hay usuario autenticado");
    }
    return u;
}
\end{lstlisting}

\subsubsection{Excepciones del dominio}

Dos excepciones propias para separar los errores de negocio de los errores de sistema:

\begin{itemize}[leftmargin=*]
  \item \texttt{NotFoundException} $\rightarrow$ HTTP 404.
  \item \texttt{ForbiddenException} $\rightarrow$ HTTP 403.
\end{itemize}

\subsubsection{Ampliación de \texttt{GlobalExceptionHandler}}

\begin{aviso}[Cambio a un archivo previo]
El \texttt{GlobalExceptionHandler} creado en la Sesión~\ref{ch:sesion3} se amplió con dos nuevos métodos:

\begin{lstlisting}[language=java]
@ExceptionHandler(NotFoundException.class)
public ResponseEntity<Map<String, Object>> handleNotFound(NotFoundException e) { ... }

@ExceptionHandler(ForbiddenException.class)
public ResponseEntity<Map<String, Object>> handleForbidden(ForbiddenException e) { ... }
\end{lstlisting}

El formato de respuesta es el mismo que ya se usaba para \texttt{IllegalArgumentException}: \texttt{timestamp}, \texttt{status}, \texttt{error}, \texttt{message}.
\end{aviso}

\subsubsection{\texttt{ChannelService}}

Contiene la lógica de canales:

\begin{itemize}[leftmargin=*]
  \item \texttt{listPublic(category, q)}: lista canales públicos con filtros opcionales.
  \item \texttt{get(id)}: devuelve un canal con el campo \texttt{subscribedByMe} calculado para el usuario actual.
  \item \texttt{create(req)}: crea un canal con el usuario autenticado como dueño.
  \item \texttt{update(id, req)}: solo permitido al dueño o a un admin.
  \item \texttt{delete(id)}: solo permitido al dueño o a un admin.
  \item \texttt{listMine()} y \texttt{listSubscribed()}: filtros por usuario.
\end{itemize}

\subsubsection{\texttt{SubscriptionService}}

La suscripción se hace con una clave compuesta (\texttt{userId}, \texttt{channelId}). El servicio verifica primero si ya existe; si existe, retorna sin error.

\subsubsection{\texttt{PostService}}

Publicación y listado. En el listado, si el canal es privado y el usuario no es dueño ni suscriptor, lanza \texttt{ForbiddenException}. El borrado solo lo puede hacer el autor o un admin.

\subsubsection{\texttt{ChannelController}}

Concentra todos los endpoints anidados bajo \texttt{/api/channels}:

\begin{lstlisting}[language=java]
GET    /api/channels
GET    /api/channels/mine
GET    /api/channels/subscribed
GET    /api/channels/{id}
POST   /api/channels
PATCH  /api/channels/{id}
DELETE /api/channels/{id}

POST   /api/channels/{id}/subscribe
DELETE /api/channels/{id}/subscribe

GET    /api/channels/{id}/posts?page=0&size=20
POST   /api/channels/{id}/posts
\end{lstlisting}

\subsection{Cambios a archivos previos}

\begin{itemize}[leftmargin=*]
  \item \texttt{GlobalExceptionHandler.java}: se añadieron dos handlers.
  \item No se modificó \texttt{SecurityConfig.java}: la regla \texttt{anyRequest().authenticated()} ya cubría \texttt{/api/channels/**}.
  \item No se modificó \texttt{application.yml}.
  \item No se modificó ninguna entidad ni migración Flyway.
\end{itemize}

\section{Verificación}

\subsection{Preparar variables}

\begin{lstlisting}[language=bash]
TOKEN=$(curl -s -X POST http://localhost:8081/api/auth/login \
  -H "Content-Type: application/json" \
  -d '{"email":"aaron@uam.mx","password":"password123"}' | jq -r .accessToken)

echo "Token: ${TOKEN:0:40}..."
\end{lstlisting}

\begin{nota}[Nota]
El error más común es ejecutar el comando de login sin \texttt{TOKEN=\$(...)} y perder el token. Sin la asignación, el token se imprime en pantalla pero no se guarda.
\end{nota}

\subsection{Crear, suscribirse y publicar}

\begin{lstlisting}[language=bash]
# Crear canal
CHANNEL_ID=$(curl -s -X POST http://localhost:8081/api/channels \
  -H "Content-Type: application/json" \
  -H "Authorization: Bearer $TOKEN" \
  -d '{"name":"Taller de Algoritmos","category":"ACADEMIC","publicChannel":true}' \
  | jq -r .id)

echo "Channel ID: $CHANNEL_ID"

# Suscribirse
curl -s -X POST "http://localhost:8081/api/channels/$CHANNEL_ID/subscribe" \
  -H "Authorization: Bearer $TOKEN" | jq

# Publicar
curl -s -X POST "http://localhost:8081/api/channels/$CHANNEL_ID/posts" \
  -H "Content-Type: application/json" \
  -H "Authorization: Bearer $TOKEN" \
  -d '{"content":"Primera publicacion"}' | jq
\end{lstlisting}

\begin{exito}[Resultado obtenido]
\begin{lstlisting}
{
  "subscribed": true,
  "channelId": "d5281a4b-8aa1-4dd4-a9e5-a86fb2425306",
  "subscriberCount": 1
}

{
  "id": "caba6493-5a2d-439c-b0ff-f589c7ff106c",
  "channelId": "d5281a4b-...",
  "channelName": "Taller de Algoritmos",
  "authorDisplayName": "Aaron Ramos",
  "content": "Primera publicacion",
  "createdAt": "2026-10-08T04:43:10.626Z"
}
\end{lstlisting}
\end{exito}

\subsection{Listar}

\begin{lstlisting}[language=bash]
curl -s "http://localhost:8081/api/channels/$CHANNEL_ID/posts?page=0&size=10" \
  -H "Authorization: Bearer $TOKEN" | jq

curl -s "http://localhost:8081/api/channels" \
  -H "Authorization: Bearer $TOKEN" | jq

curl -s "http://localhost:8081/api/channels/mine" \
  -H "Authorization: Bearer $TOKEN" | jq

curl -s "http://localhost:8081/api/channels/subscribed" \
  -H "Authorization: Bearer $TOKEN" | jq
\end{lstlisting}

\subsection{Verificar seguridad}

\begin{lstlisting}[language=bash]
curl -i -X POST http://localhost:8081/api/channels \
  -H "Content-Type: application/json" \
  -d '{"name":"Sin auth","category":"OTHER"}'
\end{lstlisting}

\begin{exito}[Resultado esperado]
\texttt{HTTP/1.1 403} sin cuerpo. La petición no trae token, por lo que Spring Security la rechaza antes de llegar al controller.
\end{exito}

\section{Resultado}

\begin{exito}[Estado al cierre]
\begin{itemize}[leftmargin=*]
  \item CRUD de canales funcional con control de acceso.
  \item Suscripciones idempotentes.
  \item Publicaciones paginadas.
  \item Filtros: por categoría, por búsqueda de texto, por canal, por usuario.
  \item Códigos HTTP correctos: 201 al crear, 204 al borrar, 403 sin token, 404 cuando no existe.
\end{itemize}
\end{exito}

\textbf{Pendiente para la siguiente sesión:} conectar la publicación de posts con Kafka para emitir eventos y alimentar el sistema de notificaciones.
